
(régulier ou non). Il a par contre le désavantage (4.6) sur @@M0@@, @@M1@@ de ne pas avoir tel quel un caractère covariant par rapport aux morphismes propres (même pour des schémas lisses quasi-projectifs sur un corps)~; il n'acquiert ce caractère qu'une fois tensorisé par @@M2@@. Autre avantage de @@M3@@ sur ses concurrents @@M4@@ et @@M5@@, c'est qu'on dispose d'un homomorphisme de @@M6@@ dans @@M7@@, permettant d'interpréter les classes de Chern cohomologiques comme images des classes de Chern à valeurs dans @@M8@@. Cet homomorphisme en général ne se factorise pas par @@M9@@ (4.7), et il est pour le moins douteux qu'il existe en général un homomorphisme de @@M10@@ dans @@M11@@ qui transforme classes de Chern en classes de Chern (Cf. n\textsuperscript{o} 5). Bien entendu, ces difficultés disparaissent lorsqu'on tensorise par @@M12@@, et on trouve un homomorphisme
@@M13@@
qui transforme classes de Chern en classes de Chern.

\textbf{4.2.} Pour établir que l'homomorphisme surjectif @@M14@@ de (4.1) est un isomorphisme modulo torsion, on utilise l'homomorphisme de Chern (déduit de la théorie des classes de Chern [9])
@@M15@@
(où le deuxième membre est l'anneau de Chern de @@M16@@ défini dans Exp.~V), qui est compatible comme on voit aisément avec la filtration topologique de @@M17@@ et la filtration naturelle de @@M18@@, et donne par passage aux gradués associés un homomorphisme de groupes gradués
@@M19@@
Ceci posé, on va prouver les formules
@@M20@@
dont la première implique d'ailleurs la seconde, grâce au fait que @@M21@@ est surjectif. Ces formules impliquent évidemment que @@M22@@ et @@M23@@ sont des isomorphismes modulo torsion, et qu'à un facteur @@M24@@ près en degré @@M25@@, ils sont inverses l'un de l'autre.

Pour prouver la première formule (4.4), on peut évidemment supposer que @@M26@@ est défini par un sous-schéma fermé irréductible @@M27@@ de @@M28@@ de codimension @@M29@@, auquel cas la formule équivaut à
